\par\medskip
\begin{center}\textls[100]{b) Die nachgelassene Abhandlung zur Theorie des agM.}\end{center}
\noindent\hspace*{\parindent}Die schier unübersehbare Fülle von neuen Resultaten und Gesichtspunkten,
die die Jahre 1799 und 1800 bei \textsc{Gauss} gezeitigt haben, mögen in
ihm den Entschluss gereift haben, einiges aus dem Schatze der erlangten Einsichten
für die Veröffentlichung zu redigieren. In der Tat macht die Abhandlung,
über die hier zu berichten ist, in den Artikeln 1—9 einen völlig
druckreifen Eindruck. Dass \textsc{Gauss} gerade mit dem elementaren Algorithmus
des agM. beginnt, ist durchaus verständlich; dass er später diese Abhandlung
abgebrochen und nicht veröffentlicht hat, werden wir nachher zu erklären
versuchen.\par
\clearpage\SourceLeaf{174}\ArithmeticHead{89}{R}{Vb. DIE NACHGELASSENE ABHANDLUNG ZUM AGM.}
\HistoricalFoot{X 2 Abh. 2.}{12}
\indent Im ersten Teil wird zunächst das agM. zwischen zwei reellen positiven
Grössen \(a\geqq b\) durch die Doppelfolge positiver Zahlen
\begingroup\makeatletter\tagsleft@true\let\veqno\leqno\makeatother
\begin{equation*}
\left\{\begin{array}{l}a,a_1,a_2,\ldots\\b,b_1,b_2,\ldots\end{array}\right.\tag{4}
\end{equation*}
\endgroup
\noindent definiert, wo die Zahlen der ersten Folge die arithmetischen, die der zweiten
die geometrischen Mittel des vorhergehenden Zahlenpaares sind; die Existenz
des gemeinsamen Grenzwerts
\[
\lim_n a_n=\lim_n b_n=M(a,b)
\]
\noindent wird streng nachgewiesen. Es folgt die Herleitung der Reihenentwicklungen
für \(M(1+k,1),\ M(1+k,1-k)\) und endlich die Reihe:
\begingroup\makeatletter\tagsleft@true\let\veqno\leqno\makeatother
\begin{equation*}
\frac1{M(1+k,1-k)}=\frac1{M(1,\sqrt{1-k^2})}
=1+(\tfrac12)^2k^2+\left(\frac{1\cdot3}{2\cdot4}\right)^2k^4+\cdots,\tag{5}
\end{equation*}
\endgroup
\noindent die uns aus älteren Aufzeichnungen bereits bekannt sind. — Im Art. 8 wird
mit Hilfe der letzteren Reihenentwicklung die lineare Differentialgleichung
zweiter Ordnung (14) Abschn. IV (S. 69) hergeleitet, der die Ausdrücke
\begingroup\makeatletter\tagsleft@true\let\veqno\leqno\makeatother
\begin{equation*}
\frac1{M(1,\sqrt{1-k^2})},\quad\frac1{M(1,k)}\tag{6}
\end{equation*}
\endgroup
\noindent Genüge leisten\footnote{Vergl. hierzu die beiden Zettel, die Werke X, 1, S. 180, III, abgedruckt sind.}, und auch gezeigt, dass
\begingroup\makeatletter\tagsleft@true\let\veqno\leqno\makeatother
\begin{equation*}
\int_0^\pi\frac{du}{(1-k^2\cos^2u)^{\frac12}}=\frac\pi{M(1,\sqrt{1-k^2})}\tag{7}
\end{equation*}
\endgroup
\noindent ist\footnote{Eine ganz direkte und sehr elegante Herleitung der Gleichung (7) gibt \textsc{Gauss} in der \textit{Determinatio attractionis} (1818, Werke III, S. 352, 353), wo er das agM. als Hilfsmittel für die Berechnung eines vollständigen elliptischen Integrals erster Gattung einführt.}. Hieraus folgt nun, dass das konstante Glied der Entwicklung von
\(\sqrt{\beta-\gamma\cos^2u}\) nach Kosinus der Vielfachen von \(u\) durch den reziproken Wert
des agM. zwischen \(\sqrt{\beta-\gamma}\) und \(\sqrt\beta\) gegeben wird. In der Scheda Ac hatte \textsc{Gauss}
diesen Zusammenhang direkt bewiesen (Werke X, 1, S. 185, 187) und daraus
durch Integration zwischen den Grenzen 0 und \(\pi\) die Gleichung (7) erschlossen.
(vergl. oben S. 67—69). Dass \textsc{Gauss} auf die Betrachtung der Entwicklungskoeffizienten
von Ausdrücken der Form \(\sqrt{\beta-\gamma\cos^2\varphi}\) durch Fragen der Störungstheorie
geführt worden ist, wurde schon (oben S. 26, 27) hervorgehoben (vergl.\par
\clearpage\SourceLeaf{175}\ArithmeticHead{90}{L}{L. SCHLESINGER, ÜBER GAUSS’ FUNKTIONENTHEORETISCHE ARBEITEN.}
\noindent den Schlusssatz des Art. 8, Werke III, S. 371 und auch einen Brief an \textsc{Olbers}
vom 3. Sept. 1805\footnote{\textit{Briefwechsel Gauss-Olbers} I, S. 269; es heisst daselbst: »Ich werde beiher meine angefangenen Arbeiten über die Störung der Planeten fortsetzen, sowohl was zur allgemeinen Theorie als zur Anwendung auf die Asteroïden gehört. In jener ist es von besonderer Wichtigkeit, die Koeffizienten, die aus
\[
\begin{aligned}
(a^2+a'^2-2aa'\cos\varphi)^{-\frac12}
&=A^0+2A'\cos\varphi+2A''\cos2\varphi\\
&\quad+2A'''\cos3\varphi+\mathrm{etc.}
\end{aligned}
\]
entspringen, leicht angeben zu können.« Diesem Briefe an \textsc{Olbers} legte \textsc{Gauss} einen Brief an \textsc{Bessel} bei, den \textsc{Olbers} an den Adressaten befördern sollte, in dem (Werke X, 1, S. 238) auch von dieser Entwicklung die Rede ist. \textsc{Gauss} schreibt daselbst: »Nun bin ich im Besitz von besonderen, zum Teil auf ganz heterogen scheinenden Untersuchungen gegründeten Kunstgriffen, jene Koeffizienten [nämlich \(A^0,A',A'',\ldots\)] zu bestimmen.« Es wird also hier auf die von \textsc{Gauss} zu jener Zeit noch geheim gehaltene Methode des agM. angespielt. Aus noch früherer Zeit wie diese beiden Briefe (wahrscheinlich Oktober 1802) dürfte eine Aufzeichnung herrühren, die sich in dem Handbuche 17, Bc, betitelt »Astronomische Untersuchungen und Rechnungen vornehmlich über die Ceres Ferdinandea (1802)« auf S. 24 in eine astronomische Untersuchung verflochten vorfindet und Werke VII, S. 384 abgedruckt ist. Vergl. hierzu die Ausführungen von M. \textsc{Brendel}, Werke XI, 2, Abh. 3, S. 216—221.}).

Für uns besonders interessant ist der Ausspruch von \textsc{Gauss} (a. a. O.
Werke III, S. 371): »Im übrigen bemerken wir, dass wir im Folgenden einen
viel allgemeineren Beweis dieser Lehrsätze geben werden, der aus viel ursprünglicheren
Prinzipien geschöpft ist; ferner werden wir alsbald lehren,
wie man auch die übrigen Koeffizienten \(Q,R,S,\ldots\) [der Entwicklung von
\((1-k^2\cos^2u)^{-\frac12}\) nach Kosinus der Vielfachen von \(u\)] durch ebenso einfache
Methoden finden kann [wie den ersten durch das agM«]. Diese Andeutung
kann sich nämlich nur auf die Verallgemeinerung der hier für die Reihe (5)
entwickelten Resultate auf Reihen allgemeineren Charakters beziehen.

\textsc{Gauss} scheint also damals, wo er eben das Gebiet der Theorie der elliptischen
Funktionen erobert hatte, die Verallgemeinerung seiner Untersuchungen
auf die Funktionen in Angriff genommen zu haben, die der jetzt sogenannten
Gaussschen Differentialgleichung genügen. Man wird dabei auch an einen
Einfluss von J. Fr. \textsc{Pfaff} zu denken haben, der in seinen \textit{Disquisitiones analyticae}
(Helmstedt 1797) in dem Kapitel »Nova disquisitio de integratione
aequationis differentio differentialis
\[
x^2(a+bx^n)ddy+x(c+ex^n)dy\,dx+(f+gx^n)y\,dx\,dx=X\,dx\,dx\text{«}
\]
\noindent diese auf \textsc{Euler} zurückgehende Differentialgleichung behandelt, und dessen
Hausgenosse \textsc{Gauss} vom Dezember 1799 bis Ostern 1800 in Helmstedt gewesen
ist (vergl. oben S. 62).\par
\clearpage\SourceLeaf{176}\ArithmeticHead{91}{R}{Vb. DIE NACHGELASSENE ABHANDLUNG ZUM AGM.}
\HistoricalFoot{}{12\textsuperscript{\(\ast\)}}
\indent In der Tat enthält die Einleitung zu der Abhandlung über die Reihe
\(F(\alpha,\beta,\gamma,x)\)\footnote{1812, Werke III, S. 123ff., vergl. Art. 6, S. 128.} die Bestimmung der Koeffizienten der Entwicklung
\[
(a^2+b^2-2ab\cos\varphi)^{-n}=A+2A'\cos\varphi+2A''\cos2\varphi+\mathrm{etc.},
\]
\noindent und auch die Methode, nach der Gauss 1800 die Differentialgleichung für
die Reihe (5) ableitet, ist dieselbe wie die, deren er sich in der \textit{Determinatio
seriei nostrae per aequationem differentialem secundi ordinis}\footnote{Nachlass, Werke III, S. 207ff.} zur Herleitung der
Differentialgleichung für die allgemeine Reihe \(F(\alpha,\beta,\gamma,x)\) bedient. Man darf
wohl auch annehmen, dass \textsc{Gauss} die Redaktion der Abhandlung von 1800
darum abgebrochen hat, weil er seinen Redaktionsplan änderte und statt von
dem Algorithmus des agM. von der allgemeinen Reihe \(F(\alpha,\beta,\gamma,x)\) ausgehen
wollte, wie er es in der Tat 1812 getan hat.

Der zweite Teil der Abhandlung von 1800 enthält nur zwei Artikel; im
ersten werden für das Differential \(dM(a,b)\) gut konvergierende Reihen angegeben,
die, wie im folgenden Artikel an Beispielen gezeigt wird, auch für
die numerische Rechnung geeignet sind. — Dies ist darum von Wichtigkeit,
weil ja die Derivierte des agM. bei der Bestimmung des Ellipsenquadranten
eine Rolle spielt, was schon in den Leistennotizen (Werke X, 1, S. 178, vergl.
oben Abschnitt IV, Gln. (3), (4)) gezeigt ist\footnote{Vergl. auch am Schluss der \textit{Determinatio attractionis}, Werke III, S. 354, Zeile 5, und am Schluss der \textit{Anzeige} ebenda, S. 360 die Formel für die Peripherie der Ellipse.}. Besonders die Reihe
\[
\frac12+\frac14\frac{a-a_1}{a_1}+\frac18\frac{a-a_1}{a_1}\cdot\frac{a_1-a_2}{a_2}+\cdots,
\]
\noindent wo die \(a,a_1,a_2,\ldots\) die Elemente der ersten Folge (4), also die arithmetischen
Mittel bedeuten, konvergiert sehr stark; sie findet sich auch in der Werke
X, 1, S. 207, VI, abgedruckten Notiz der Scheda Ae aus dem Juli 1800 mit
einigen weiteren Ausführungen und wird uns später noch in der Korrespondenz
von \textsc{Gauss} mit \textsc{Schumacher} (1816) begegnen.

Was nun die Frage der Bestimmung von \(t=\frac{\varpi'}\varpi\) für einen gegebenen
Wert von \(\mu\) anlangt, so sehen wir, dass für \(k=\sin v\) die beiden partikularen
Lösungen (6) der von \textsc{Gauss} aufgestellten linearen Differentialgleichung gerade
die zur Herstellung des Quotienten (2\,a) erforderlichen Elemente liefern. Die
Entwicklung (5) von \(\frac1{M(1,\sqrt{1-k^2})}\) nach positiven ganzen Potenzen von \(k^2\) war\par
\clearpage\SourceLeaf{177}\ArithmeticHead{92}{L}{L. SCHLESINGER, ÜBER GAUSS’ FUNKTIONENTHEORETISCHE ARBEITEN.}
\noindent \textsc{Gauss} geläufig, die Darstellung von \(\frac1{M(1,k)}\) in der Umgebung von \(k=0\) ergibt
sich entweder aus der Differentialgleichung (14) Abschn. IV (oben S. 69)
oder direkt aus der Formel
\[
M(x,1)=\frac{C(x-\alpha x^{-1}-\beta x^{-3}-\gamma x^{-5}-\cdots)}{\log(4x-ax^{-1}-bx^{-3}-cx^{-5}-\cdots)}
\]
\noindent der Scheda Ac (Werke X, 1, S. 186, vergl. oben Gl. (10) Abschn. IV) als
\begingroup\makeatletter\tagsleft@true\let\veqno\leqno\makeatother
\begin{equation*}
\frac1{M(1,k)}=\frac1\pi\frac1{M(1,\sqrt{1-k^2})}
\left(\log\frac{16}{k^2}-\frac{k^2}2+\delta_1k^4+\cdots\right).\tag{8}
\end{equation*}
\endgroup
\noindent Diese Darstellung konvergiert ebenso wie die Reihe (5) selbst für \(|k^2|<1\). —
Nun kommt in der Scheda Ac (Werke X, 1, S. 198) die Transformation
\begingroup\makeatletter\tagsleft@true\let\veqno\leqno\makeatother
\begin{equation*}
\sin\psi=\operatorname{tg}^2\tfrac12v\tag{9}
\end{equation*}
\endgroup
\noindent vor, deren Bedeutung darin besteht, dass sie für \(\sin\psi=k_1\),
\begingroup\makeatletter\tagsleft@true\let\veqno\leqno\makeatother
\begin{equation*}
k_1=\frac{1-k'}{1+k'},\quad k'=\sqrt{1-k^2}\tag{9\,a}
\end{equation*}
\endgroup
\noindent ergibt, also auf den Algorithmus des agM. hinauskommt, indem, wenn man
\(a=1,b=k',c=\sqrt{1-k'^2}=k\) nimmt, und hieraus nach dem Algorithmus
des agM.
\[
a_1=\frac{1+k'}2,\quad b_1=\sqrt{k'},\quad c_1=\sqrt{a_1^2-b_1^2}=\frac{1-k'}2
\]
\noindent bildet,
\[
\frac{b_1}{a_1}=k'_1=\sqrt{1-k_1^2},\quad\frac{c_1}{a_1}=k_1
\]
\noindent ist. Schreibt man nun für die reziproken Werte der \(M(1,k'_1),M(1,k_1)\) die
den Entwicklungen (5) und (8) analogen Entwicklungen nach Potenzen von
\(k_1^2=\nu_1\) hin, so konvergieren diese für \(|\nu_1|<1\), also vermöge der Relation
(9\,a) für alle komplexen Werte von \(k^2=\nu\). Wegen \(M(a,c)=2M(a_1,c_1)\) ist
aber
\[
t_1=\frac{M(1,\cos\psi)}{M(1,\sin\psi)}=2\frac{M(1,\cos v)}{M(1,\sin v)}=2t,
\]
\noindent also mit \(t_1\) auch \(t\) für alle Werte von \(k^2\) definiert. — Dass \textsc{Gauss} diese Anwendung
der Transformation (9) im Auge hatte, zeigt der Art. 5 (Werke X, 1,
S. 197—198) der Scheda Ac-Aufzeichnung, der sonst gegenüber dem Art. 7
der Aufzeichnung IV (Werke X, 1, S. 190) nichts wesentlich Neues enthält.
— Durch wiederholte Anwendung der Transformation (9) kann man für die\par
\clearpage\SourceLeaf{178}\ArithmeticHead{93}{R}{Vc. THEORIE DER MODULFUNKTION.}
\noindent agM. immer besser konvergierende Entwicklungen erzielen, ein Umstand, den
z. B. auch \textsc{Weierstrass}\footnote{Sitzungsberichte der Berliner Akademie 1883, \textsc{Weierstrass'} Werke III, S. 257ff. Vergleiche H. \textsc{Weber} a. a. O.\textsuperscript{102)}, S. 144.} ausgenutzt hat.
